% 02_setup.tex - unpack the workshop, set the two variables, install the Python packages
\section{Set up}

\step{Unpack the workshop}\label{st:unpack}
Copy \cmd{Workshop\_1\_Getting\_Started.zip} into your home folder. On WSL, your Windows files are under \cmd{/mnt/c/}, for
example:
\begin{shell}
$ cp /mnt/c/Users/YOUR_NAME/Downloads/Workshop_1_Getting_Started.zip ~/
\end{shell}
Unpack the zip, go into the new folder, and set two variables. \cmd{WS} is the workshop folder and \cmd{RAVEN} is the
Raven program:
\begin{shell}
$ cd ~
$ unzip -q Workshop_1_Getting_Started.zip
$ cd Workshop_1_Getting_Started
$ export WS=$PWD
$ export RAVEN=$WS/bin/Raven
\end{shell}
\see no output. \cmd{unzip -q} and \cmd{export} print nothing when they work. Now list the folder:
\term{t1_ls_ws.txt}
\see these names (your terminal may show them in another order or arrangement).
\cmd{bin} holds Raven and the MODFLOW~6 library, \cmd{model/start} the Liard River example as you get it,
\cmd{model/finished} the one file that this workshop changes (\cmd{Liard.rvi}; its \cmd{README.txt} explains), \cmd{scripts} the helper programs, and
\cmd{tex} the source of this document.
\begin{trouble}
\begin{itemize}
\item \cmd{unzip: command not found}: install it with \cmd{sudo apt install unzip} (Ubuntu and WSL).
\item \cmd{cannot find or open Workshop\_1\_Getting\_Started.zip}: the zip is not in your home folder. Type \cmd{ls \textasciitilde}
  to see what is there.
\end{itemize}
\end{trouble}
\begin{note}[Every new terminal]
\cmd{export} is valid only in the terminal where you typed it. In a new terminal, go to the workshop folder and type the
two \cmd{export} lines again.
\end{note}

\step{Install the Python packages}\label{st:python}
The helper scripts need Python~3 with numpy, pandas and matplotlib. A \emph{virtual environment} (a private folder of
Python packages) avoids conflicts with the packages of your system. On Ubuntu and WSL:
\begin{shell}
$ sudo apt update
$ sudo apt install python3-venv
$ python3 -m venv ~/raven-venv
$ source ~/raven-venv/bin/activate
$ pip install numpy pandas matplotlib
\end{shell}
\see \cmd{pip} downloads the packages. It prints a line that starts with \cmd{Successfully installed}. Your prompt now starts with
\cmd{(raven-venv)}. Test the packages:
\begin{shell}
$ python3 -c "import numpy, pandas, matplotlib; print('ok')"
ok
\end{shell}
\begin{note}[Every new terminal]
Type \cmd{source \textasciitilde/raven-venv/bin/activate} in every new terminal, before you run a Python script.
\end{note}
